\documentclass[UTF8,11pt,a4paper,fontset=fandol]{ctexart}
\usepackage[margin=20mm]{geometry}
\usepackage{amsmath,amssymb,booktabs,array,tabularx,enumitem,setspace,xcolor,hyperref,xurl,fancyhdr}
\setstretch{1.08}
\setlength{\emergencystretch}{3em}
\setlength{\headheight}{15pt}
\setlength{\tabcolsep}{5pt}
\renewcommand{\arraystretch}{1.2}
\setlist{nosep,leftmargin=2em}
\hypersetup{colorlinks=true,linkcolor=blue!45!black,urlcolor=blue!45!black}
\providecommand{\WorkloadName}{工作负载}
\pagestyle{fancy}\fancyhf{}
\fancyhead[L]{\small Table II(b)：算子级 RI}
\fancyhead[R]{\small\WorkloadName}
\fancyfoot[C]{\thepage}
\newcommand{\qs}{Q_{\mathrm S}}
\newcommand{\qr}{Q_{\mathrm R}}
\newcommand{\RI}{\mathrm{RI}}
\newcommand{\para}[1]{\par\smallskip\noindent\textbf{#1}\quad}
\newcommand{\ModelTableSetup}{\renewcommand{\arraystretch}{1.2}\setlength{\tabcolsep}{4pt}}
\newcommand{\ResultTableSetup}{\renewcommand{\arraystretch}{1.45}\setlength{\tabcolsep}{4pt}}
\newcommand{\DisplayNote}{\par\smallskip{\noindent\footnotesize $1\mathrm{K}=1024$，$1\mathrm{M}=1024^2$。RI 至少为 1 时保留一位小数，小于 1 时保留三位有效数字；精确值见机器数据。\par}\smallskip}
\author{}\date{2026 年 10 月 2 日}

% Compact, identical title block keeps the mathematical introduction on one page.
\makeatletter
\renewcommand{\maketitle}{\thispagestyle{plain}\begin{center}
{\LARGE\bfseries\@title\par}\vspace{8pt}{\normalsize\@date\par}
\end{center}\vspace{3pt}}
\makeatother

\renewcommand{\WorkloadName}{注意力}
\title{\bfseries\WorkloadName 的输入复用与 RI}
\begin{document}
\maketitle
\section{对象与公式}
取单序列、一个完整因果 GQA 层的 QK 与 AV 两个矩阵阶段。
$H_q$ 为 query 头数，$H_{\rm kv}$ 为 KV 头数，$g=H_q/H_{\rm kv}$；
$d_{\rm QK}$ 与 $d_V$ 为原生头维，$L$ 为本次追加后可见长度。
每个 KV head 保存一份 K/V，由组内 $g$ 个 query heads 共享。

沿用 Table II(a) 的 $\RI=\qs/\qr$ 与每元素 1 Byte 口径。
在有效前缀 $i$，K 的形状为 $i\times d_{\rm QK}$，V 为 $d_V\times i$；
每个 query head 在 QK 阶段输入 $d_{\rm QK}$ 个 query 元素，
在 AV 阶段输入 softmax 后的 $i$ 个系数。历史 K/V 是 resident 操作数，不重复列为 streaming 输入。
两阶段输入分别计入，输出不另加；
$\qr$ 只计窗口内新写入的唯一逻辑 K/V。

\para{预填充（Prefill）}从空 KV 建到 $L$，计完 $L$ 个位置的因果有效范围：
\[
\begin{gathered}
\qs=\sum_{i=1}^{L}H_q(d_{\rm QK}+i)=H_q\!\left[Ld_{\rm QK}+\frac{L(L+1)}2\right],
\qquad \qr=LH_{\rm kv}(d_{\rm QK}+d_V),\\[4pt]
\boxed{\displaystyle\RI=\frac{g[d_{\rm QK}+(L+1)/2]}{d_{\rm QK}+d_V}.}
\end{gathered}
\]
这是三角求和：第 $i$ 个 query 使用当时可见的 $i$ 个 K/V，AV 输入随前缀增长。
每个 token 的 K/V 只写一次，故写入是 $L$ 份新状态之和，不是每次重写整个前缀。

\para{单步解码（Decode）}已有 $L-1$ 个 token 的 KV，写入一份新的 K/V 后，对长度 $L$ 的完整前缀求值：
\[
\begin{gathered}
\qs=H_q(d_{\rm QK}+L),\qquad \qr=H_{\rm kv}(d_{\rm QK}+d_V),\\[4pt]
\boxed{\displaystyle\RI=\frac{g(d_{\rm QK}+L)}{d_{\rm QK}+d_V}.}
\end{gathered}
\]
最终有效 KV 容量均为 $LH_{\rm kv}(d_{\rm QK}+d_V)$ Byte；
初始容量在 Prefill 为 0，在 Decode 为 $(L-1)H_{\rm kv}(d_{\rm QK}+d_V)$ Byte。
容量与本窗口的累计写入量分别记录。

\section{模型信息}
六模型均取固定版本中的完整/global GQA 层，层号从 0 起。
\begin{center}\small\ModelTableSetup
\begin{tabularx}{\linewidth}{@{}p{0.32\linewidth}*{6}{>{\centering\arraybackslash}X}@{}}
\toprule
模型 & 所选层 & $H_q$ & $H_{\rm kv}$ & $g$ & $d_{\rm QK}$ & $d_V$\\
\midrule
Qwen3.5-2B & 3 & 8 & 2 & 4 & 256 & 256\\
Ministral 3 8B (2512) & 0 & 32 & 8 & 4 & 128 & 128\\
Qwen3.6-35B-A3B & 3 & 16 & 2 & 8 & 256 & 256\\
Tencent Hy3 (295B) & 1 & 64 & 8 & 8 & 128 & 128\\
Ling-1T & 4 & 64 & 8 & 8 & 128 & 128\\
MiMo-V2.5-Pro & 7 & 128 & 8 & 16 & 192 & 128\\
\bottomrule
\end{tabularx}
\end{center}
各模型保留原有上下文设置；Ling-1T 的 64K 沿用既定官方扩展配置。\footnote{YaRN：factor=4、original\_max\_position\_embeddings=32768、type=yarn，并设 \texttt{--max-model-len 131072}；原始配置不改。}
MiMo 的 Value 在写入前乘 0.612，维度不变；所选全局层无 sink。

\newpage
\section{代入结果}
两表以模型为行、可见长度 $L\in\{1\mathrm{K},8\mathrm{K},64\mathrm{K}\}$ 为列，单元格为小数 $\RI$。
\begin{center}\small\textbf{预填充（Prefill）的 $\RI$：空 KV 建到 $L$}\par\smallskip
\ResultTableSetup
\begin{tabularx}{\linewidth}{@{}p{0.32\linewidth}*{3}{>{\centering\arraybackslash}X}@{}}
\toprule
模型 & $L=1\mathrm{K}$ & $L=8\mathrm{K}$ & $L=64\mathrm{K}$\\
\midrule
Qwen3.5-2B & $6.0$ & $34.0$ & $258.0$\\
Ministral 3 8B (2512) & $10.0$ & $66.0$ & $514.0$\\
Qwen3.6-35B-A3B & $12.0$ & $68.0$ & $516.0$\\
Tencent Hy3 (295B) & $20.0$ & $132.0$ & $1028.0$\\
Ling-1T & $20.0$ & $132.0$ & $1028.0$\\
MiMo-V2.5-Pro & $35.2$ & $214.4$ & $1648.0$\\
\bottomrule
\end{tabularx}
\par\bigskip
\textbf{单步解码（Decode）的 $\RI$：已有 $L-1$，追加一个}\par\smallskip
\ResultTableSetup
\begin{tabularx}{\linewidth}{@{}p{0.32\linewidth}*{3}{>{\centering\arraybackslash}X}@{}}
\toprule
模型 & $L=1\mathrm{K}$ & $L=8\mathrm{K}$ & $L=64\mathrm{K}$\\
\midrule
Qwen3.5-2B & $10.0$ & $66.0$ & $514.0$\\
Ministral 3 8B (2512) & $18.0$ & $130.0$ & $1026.0$\\
Qwen3.6-35B-A3B & $20.0$ & $132.0$ & $1028.0$\\
Tencent Hy3 (295B) & $36.0$ & $260.0$ & $2052.0$\\
Ling-1T & $36.0$ & $260.0$ & $2052.0$\\
MiMo-V2.5-Pro & $60.8$ & $419.2$ & $3286.4$\\
\bottomrule
\end{tabularx}
\end{center}
\DisplayNote
例如，Qwen3.5-2B 在 Prefill、$L=1\mathrm{K}$ 时，
$\qs=8[1024\times256+1024\times1025/2]=6\,295\,552$ Byte，
$\qr=1024\times2\times(256+256)=1\,048\,576$ Byte，
故 $\RI\approx6.0$。

\section{结果说明}
固定头维时，RI 随 $g$ 成比例增长。两个 Qwen 模型的头维相同，
Qwen3.6 的 $g$ 加倍，RI 也加倍。

Prefill 累计全部 $L$ 个位置，Decode 只计算当前位置；后者的写入仅为前者的 $1/L$。
同一 $L$ 下，$\RI_{\rm Decode}/\RI_{\rm Prefill}=(d_{\rm QK}+L)/[d_{\rm QK}+(L+1)/2]$，
长前缀时趋近 2。该关系来自当前窗口和因果前缀计数。

Hy3 与 Ling 的头结构相同，RI 相同；MiMo 直接使用原生 $192/128$ 头维。
\par\smallskip\noindent\small
固定模型来源、精确字节及复算记录见本目录 README 与 data/；数学入口为 Table II(a)。
\end{document}
